\section{Setting and conventions}
\label{sec:setting}

\subsection{The Hamiltonian}

Work on the cubic lattice $\Z^3$ with gauge group $G = \SU(N)$, one
group element $U_\ell$ per oriented link. The Kogut--Susskind
Hamiltonian is
\begin{equation}
  H \;=\; H_0 + V,
  \qquad
  H_0 = \tfrac12 \sum_{\ell} E_\ell^2,
  \qquad
  V = -\,u\, W,
  \label{eq:KS}
\end{equation}
where $E_\ell^2$ is the quadratic Casimir of the left-invariant electric
field on link $\ell$, and $W$ is the sum of plaquette characters in both
orientations,
\begin{equation}
  W = \sum_{p} \big(\chi_p + \bar\chi_p\big),
  \qquad
  \chi_p = \Tr \prod_{\ell \in \partial p} U_\ell .
\end{equation}

Strong-coupling perturbation theory is the expansion in $u$ around the
electric vacuum, organized as degenerate perturbation theory --- a des
Cloizeaux effective operator --- on the sector of interest.

\subsection{The coupling, and one erratum that matters}

The canonical coupling variable throughout is
\begin{equation}
  u \;=\; \frac{\beta_N}{2N},
  \label{eq:coupling}
\end{equation}
whose $\SU(3)$ specialization is $u = \beta_3/6$.

A warning, because it has corrupted transcriptions of this material
before. An archived line reading $Y = 2\beta/3 = 4u$ is a
definition-label erratum, not a change of variable. The coefficients
printed alongside it were \emph{already} in the canonical coordinate
\eqref{eq:coupling}. Applying a $4^r$ rescaling at order $r$ to
``correct'' them corrupts every order, and the repository carries an
explicit check that the printed towers reproduce their canonical values
verbatim.

For rank-uniform statements the natural variable is the planar one,
\begin{equation}
  \tau \;=\; \frac{2u}{N^2} \;=\; \frac{\beta_N}{N^3},
  \label{eq:tau}
\end{equation}
in which the grading theorems of \S\ref{sec:grading} take their simplest
form: every even order of the one-plaquette splitting then carries the
same single power of $N$.

\subsection{Two projections that must not be conflated}

Two distinct projections appear throughout, and several contradictions
in this program's own register were naming collisions between them.

The \emph{band} projection is onto the degenerate strong-coupling
subspace carrying the excitation of interest --- for the flux band, the
cellular subspace of \S\ref{sec:carrier}. Coefficients anchored to it
are band-kernel quantities.

The \emph{physical} projection is onto the gauge-invariant,
vacuum-subtracted spectrum at a definite momentum. Coefficients anchored
to it are physical quantities.

These are related by translation-local scalar shifts, and such a shift
changes neither the centered operator, nor its eigenvectors, nor the
bandwidth. It does change the scalar. Consequently the band-kernel
anchor and the physical $\Gamma$-point coefficient are two differently
anchored coordinates, not two competing estimates of one quantity, and
they are written $q^{(4)}_{\mathrm{band}}$ and $m^{(4)}_{\Gamma}$
respectively. Writing either as ``$m_4$'' regenerates a contradiction
that does not exist, and the repository's search tooling refuses the
name for that reason.

\subsection{Exactness}

Every coefficient displayed in Part~\ref{part:exact} is an exact
rational, or an exact rational function of $N$. Where a quantity is
known here only numerically it carries an explicit numerical marker in
the source registry and is never displayed as though exact. The single
most dangerous defect available in this material is a float that reads
as exact, and the boundary between the two is machine-guarded rather
than maintained by care.
